% chapters/15/section-05.tex
\section{Partikel Sederhana}
\label{sec:partikel}

\subsection{Pasangan yang Terlihat bagi Bunyi}

Bagian ini bersifat pilihan menurut RPS, tetapi ia melengkapi gagasan yang
dibahas sepanjang bab ini. Setiap bunyi sebaiknya memiliki pasangan yang
terlihat. Untuk beberapa peristiwa, pasangannya sudah ada: skor bertambah
ketika benda diambil, nyawa berkurang ketika pemain terluka. Untuk peristiwa
lain, terutama yang berlangsung sekejap, pasangannya belum ada. Di situlah
partikel berguna.

Partikel adalah kumpulan gambar kecil yang muncul, bergerak, lalu menghilang.
Satu ledakan kecil, satu kepulan debu ketika pemain mendarat, atau satu kilau
ketika benda diambil. Ketiganya berlangsung kurang dari satu detik, dan
ketiganya memberi pemain bukti yang terlihat bahwa sesuatu benar-benar terjadi.

\subsection{Memasang Sistem Partikel}

Komponen \istilah{Particle System} diambil dari menu \istilah{Component}, lalu
\istilah{Effects}, lalu \istilah{Particle System}. Susunan bawaannya sudah
memperlihatkan sekumpulan partikel yang menyembur ke atas. Empat pengaturan
berikut yang paling perlu Anda sesuaikan.

\begin{description}
  \item[\istilah{Duration} dan \istilah{Looping}] Lama satu putaran partikel
        dan apakah ia diulang. Untuk bunyi yang dipicu peristiwa, nilai
        \istilah{Looping} dimatikan, sebab partikelnya diputar sekali saja.
  \item[\istilah{Start Lifetime}] Lama setiap partikel hidup sebelum hilang.
        Untuk percikan kecil, nilai seperlima sampai setengah detik sudah
        cukup.
  \item[\istilah{Start Speed}] Kecepatan awal partikel meninggalkan pusatnya.
        Nilai besar memberi kesan ledakan; nilai kecil memberi kesan kepulan.
  \item[\istilah{Emission}] Jumlah partikel yang dikeluarkan setiap satuan
        waktu. \istilah{Bursts} dipakai untuk mengeluarkan sekumpulan partikel
        sekaligus pada satu saat, dan itulah yang sesuai untuk peristiwa
        sekejap.
\end{description}

\subsection{Memanggil Partikel pada Peristiwa yang Sama}

Partikel dipanggil pada peristiwa yang sama dengan bunyinya, sehingga
keduanya muncul bersamaan. Caranya melalui \code{ParticleSystem.Play()},
sebagaimana bunyi dipanggil melalui \code{AudioSource}.

\begin{kodekeluaran}
// Ditambahkan pada PemicuBunyi, disimpan pada peubah partikel.
public ParticleSystem partikelAmbil;

void OnTriggerEnter2D(Collider2D lain)
{
    if (!lain.CompareTag("Pickup")) return;

    if (manajer != null) manajer.BunyiAmbil();
    if (partikelAmbil != null) partikelAmbil.Play();
    ...
}
\end{kodekeluaran}

Partikelnya dipanggil pada tempat yang sama persis dengan
bunyinya. Itu bukan kebetulan, melainkan akibat dari cara kerja yang sama:
keduanya adalah tanggapan atas satu peristiwa, dan keduanya sebaiknya
diurus pada satu tempat agar tidak ada yang tertinggal.

\begin{peringatanbox}
Jangan memakai partikel untuk menutupi masalah. Apabila sebuah mekanik tidak
terasa jelas tanpa partikel, kemungkinan besar mekaniknya sendiri belum jelas.
Partikel ditambahkan setelah mekaniknya bekerja dan terasa benar, bukan sebagai
tambalan atas mekanik yang membingungkan.
\end{peringatanbox}
